
PRINCIPLES

E. E. Gynild\footnote{Gynild was president of the Augsburg Seminary Board of Trustees in 1918. Ragna Sverdrup served as treasurer.}


Installment 1: “Folkebladet,” March 20, 1918

Principles are an excellent thing to possess. A man without principle is like a ship without a rudder. To imagine a Christian without principle is surely an impossibility. To walk in Jesus’ footsteps is surely to have a principle? Wherever a Christian becomes indifferent in his following of his Savior, there the Christian power is broken and has lost the ability to influence others in a Christian direction.

From the moment a person is saved, that same person is given a new direction of life. The imparting of the Holy Spirit and the love of Christ poured out in his heart effects this. The world has never stood face to face with a stronger life-principle than that which it met in people sealed by the Spirit of God. Naturally, neither has the world been able to prevent the progress of Christianity. And yet the life of the Christians was a life in freedom. Yes, the only life that is free. The person of worldly mind becomes nothing but a weather vane that turns with the wind. Such was not the way of the Lord’s witnesses. “What did you go out into the wilderness to see? A reed moved back and forth by the wind?” No, they will not find that in John. Neither king nor scaffold forces him away from his straight course. The Forerunner goes to death, even as did he for whom he prepared the way.

And yet people venture to speak scornfully of principles. But the world has never liked people with Christian principles. All of them have been, more or less, intolerable to it, and hosts of them have laid down their lives for their testimony and received the martyr’s crown. But when they lay in their graves, many joined in wreathing those graves.

If there is anything our age needs, it is people with Christian principles, people for whom “truth is life’s highest goal,” truth in life and in work.

That not everyone shares the same view, either in state or church, is of course a well-known fact. But though others may hold another view concerning Christian matters, and still more concerning ecclesiastical matters, that will give us no excuse if we frivolously play with the truth. We must will truth. Only then does our life become whole. The Jesuitical will then be driven out.

When church people today maintain that among the various Lutheran church bodies there is no difference in viewpoint and manner of work, those people are either ignorant, or they are full of deceit.

Within the Norwegian Lutheran Church there have, during the last century, been two ecclesiastical tendencies: a high-church and a low-church tendency, a state-church and a congregational tendency, even though the latter has not always been conscious of what it was striving toward. The same tendencies are still present in the Norwegian Lutheran Church on both sides of the ocean. As time has advanced, there have been strong tendencies and a more intense striving to erase every difference of principle and to lead all tendencies into one united great stream called the Church. It is to be made great and broad so that everything may find room. The church body is to make room for everything. But this is something that cannot possibly be done. For if principles and tendencies are to be allowed to live, they must of necessity come into conflict; and if they do not possess within themselves enough power of truth to enter into conflict, then they do not deserve the name of tendencies or Christian principles. But here the stifling element enters. The Church, the church body, wants peace. Everything is to be sacrificed for the church body; and what then becomes of the tendencies? The Church, this all-devouring thing, which was to be a means for advancing the work of the kingdom of God, has become the goal; the majority of the church body has become the whole direction. The demand that everything be sacrificed for the church body is a very strong poison against freedom, both individually and congregationally. It is the old Catholic principle in a somewhat diluted degree. This kind of weapon has proved by no means ineffective. For people so dearly want a church without conflict. But when one will not fight, it always goes badly. He who lays down his weapons must share Russia’s fate, as we see it today.\footnote{An allusion to the Treaty of Brest-Litovsk, signed March 3, 1918, by which Soviet Russia withdrew from the First World War on severe terms.} This voice, saying that everything is alike, comes from the same abyss and has the same purpose: to kill the various tendencies and render principles worthless.

You who now believe that the free congregation is worth our labor and sacrifices, and that in this we serve the Lord’s cause, give now to Augsburg what you are able to sacrifice, so that the debt may be lifted off!

(To be continued.)

PRINCIPLES
(E. E. Gynild.)

Installment 2: “Folkebladet,” March 27, 1918

(Conclusion.)

It is not known to me that anyone has yet attempted to prove that the Guiding Principles of the Lutheran Free Church are not in full agreement with the New Testament revelation concerning the congregation. In spite of this, there have been enough who have censured the Lutheran Free Church. But then the censure must chiefly fall upon us who constitute it and labor within it, because we have not been equal to carrying the principles through.

But if those same principles demanded no more than what we were able to carry through, then they would not be Christian principles. But if the charge concerns us as congregational people, as workers who ought willingly to present ourselves as instruments of the Spirit, then we are not blameless, neither laypeople nor pastors.

It has pleased the Lord to use his people for the accomplishment of the work in and through his congregation. Here we are sluggish. That rebuke we must receive to our humbling.

When the broad-church tendency comes after us, it is chiefly because, they say, we are like them. And to that we have nothing to say when they look upon us as individuals. But when one says that we seek the same thing and speaks of us as the Free Church, then the matter looks somewhat different. Yet here is an opportunity for them to prove that it is so. They can simply adopt the Guiding Principles of the Lutheran Free Church and, if they wish, the constitution that has sprung from them, and then they have given the proof.

The other tendency accuses us because we do not use our congregations in harmony with our principles. To them we answer that we believe we are in harmony with the principles. They must surely mean that we have not reached the goal, and in that they are entirely right. \textbf{A “free and living” congregation is not attained in the sense that we can say, Now we have it. That would then have to be a human work.} The longing to see a living and liberated congregation surely lives in every believing heart, but the goal is so high that it will occupy all our time and all our powers and gifts, and even then the goal lies far ahead. If we were at the goal, there would surely be nothing left to expect but decline. But when we take up the work for the “free and living congregation,” then we shall not outlive our calling in life.

Some have thought that one should attempt to tear out of the existing congregational organizations what, in our judgment, constitutes suitable material for a “living congregation.” In the first place, we human beings are not altogether infallible judges, so that we would be fit for such a separating-out; and in the second place, our principles forbid us this kind of procedure.

“The making alive and setting free of the congregations” is what lies before us in the Lutheran Free Church. It is the work of the Holy Spirit. As workers, our part is to labor and not grow weary. That there will continually be some who grow weary, who for one or another carnal reason fail, will unfortunately no doubt continue to be the case. But such falls also we must try to see in Christian light and not through the spectacles of church politics. The cause lies deeper than an outward connection with this or that. He who betrays his ideals fares badly and deals ill with himself. Yes, it is best not to say more about so unpleasant a subject.

Much has lately been written about liberation, but little about life. Yet life is the first condition of liberation. The Apostle Peter joins the two in this way: “Come to him, the living Stone”—and yourselves, as living stones, “be built into a spiritual house, a holy priesthood,” and so forth. He who would labor for the liberated congregation must heed the whole counsel of God. It will not do to take only that which sounds like us, or which for the sake of one purpose or another is expedient. A sled has two runners, and the load is carried forward when both runners bear the load. But perhaps we do as the advocate of dead orthodoxy does: we presuppose the life? Then we shall also obtain the same result: “death in the pot.”

In the free congregation no one has the right to claim unconditional authority, unless he can say, as the prophets did, ``Thus says the Lord.''\footnote{The printed text reads, somewhat corruptly, ``Ikke uten ham. Kan sige, som profeterne.'' The sense is supplied by the immediate context.} Our authority is our Bible. He who wishes to rule, and has gifts and powers for doing so, must subdue and kill too much around him, whether the ruler is the pastor in “his” congregation or a brother among brothers. Many will indeed remain silent before superior force, but that kind of resignation bears as its fruit a weak will, a dangerous spiritless condition, which is of little use in the work for a “free and living congregation.”

Now, a word to all of us who have taken our stand upon the principles of the Lutheran Free Church. Have we done what we could to carry them forward toward a clearer breakthrough among our people? Perhaps we have become lukewarm, so that they no longer have the place in our life and work that they once had? If they still possess our love, then let us gather around Augsburg, the school that has nurtured them and continually carried them forward among our people. The work for the congregation is heavy enough without our also bearing an unnecessary burden, an oppressive debt, which is already too old to be tolerated any longer. One more concerted effort is needed, and we must make it now, in this year!
